\subsection{基本定义与假设}\label{sec:conservation-basic-definitions}

首先计算守恒律的特征, 并给出使方程组严格双曲的条件, 由此开始对守恒律的研究.

\begin{Lemma}
\label{lem:conservation-characteristics}
若矩阵
\begin{equation}\label{eq:conservation-characteristic-matrix}
\hat x'(t)\bm{I}-\nabla\bm{f}(\bm{u}(\hat x(t),t))
\end{equation}
奇异, 则曲线 $t\mapsto\hat x(t)$ 是具有解 $\bm{u}(x,t)$ 的守恒律 \eqref{eq:conservation-law-system} 的特征曲线. 此外, 若 $\nabla\bm{f}(\bm{u})$ 的特征值为互不相同的实数, 则方程组在解 $\bm{u}$ 处严格双曲.
\end{Lemma}

证明留给读者. 所需的全部工作, 就是在用函数图像而非水平集描述曲线的情形下, 解释非线性方程组的特征曲线与严格双曲性的定义. (回忆 Section~\ref{sec:cauchy-kovalevskaya} 中关于曲面的不同表示方式的说明.)

当然, 特征曲线的斜率就是矩阵 $\nabla\bm{f}(\bm{u})$ 的特征值. 据此引入描述 $\nabla\bm{f}$ 的特征值和特征向量的一些记号. 假设方程组严格双曲, 从而有 $n$ 个互不相同的实特征值 $\lambda_1(\bm{u})<\cdots<\lambda_n(\bm{u})$, 对应的右特征向量 $\bm{r}_k(\bm{u})$ 和左特征向量 $\bm{l}_k(\bm{u})$ 满足
\begin{equation}\label{eq:conservation-right-eigenvectors}
\nabla\bm{f}(\bm{u})\bm{r}_k(\bm{u})=\lambda_k(\bm{u})\bm{r}_k(\bm{u}),
\end{equation}

\begin{equation}\label{eq:conservation-left-eigenvectors}
\bm{l}_k(\bm{u})^\mathrm{T}\nabla\bm{f}(\bm{u})=\lambda_k(\bm{u})\bm{l}_k(\bm{u})^\mathrm{T}.
\end{equation}
回忆, 由于特征向量互不相同, 右特征向量集 $\{\bm{r}_1(\bm{u}),\allowbreak \ldots,\bm{r}_n(\bm{u})\}$ 与左特征向量集 $\{\bm{l}_1(\bm{u}),\allowbreak \ldots,\bm{l}_n(\bm{u})\}$ 各自构成状态空间的一组基.

现在在状态空间上定义一些称为 \emph{Riemann 不变量}的函数, 它们在求解具有不连续初始条件的问题时十分重要. 这些函数局部定义在一个邻域 $U\subseteq\mathbb{R}^n$ 上.

\begin{Definition}
\label{def:riemann-invariant}
$k$-Riemann 不变量是一个光滑函数 $w:U\to\mathbb{R}$, 它对每个 $\bm{u}\in U$ 满足
\begin{equation}\label{eq:riemann-invariant-definition}
\bm{r}_k(\bm{u})\cdot\nabla w(\bm{u})=0.
\end{equation}

\end{Definition}

以下引理给出一个适当的 Riemann 不变量组的存在性结果.

\begin{Lemma}
\label{lem:riemann-invariants-existence}
对每个 $\bar{\bm{u}}\in\mathbb{R}^n$, 都存在 $\bar{\bm{u}}$ 的邻域 $\bar U\subset\mathbb{R}^n$, 其上有 $n-1$ 个 $k$-Riemann 不变量, 它们的梯度在每个点 $\bm{u}\in\bar U$ 处线性无关.
\end{Lemma}

\begin{Proof}
设 $S$ 为通过点 $\bar{\bm{u}}$ 并与向量 $\bm{r}_k(\bar{\bm{u}})$ 横截的光滑曲面. 在 $\bar{\bm{u}}$ 的一个邻域内, 考虑常微分方程组 $\mathrm{d}\bm{u}/\mathrm{d}t=\bm{r}_k(\bm{u})$. 若 $w(\bm{u})$ 沿这一常微分方程组的每条轨线为常数, 它就是一个 $k$-Riemann 不变量. 现在, 每条经过 $\bar{\bm{u}}$ 的充分小邻域的轨线都恰好与 $S$ 相交一次. 交点的坐标 (相对于 $S$ 上适当选定的坐标系) 就作为所需的 Riemann 不变量.
\end{Proof}

\begin{Example}{$p$ 方程组}
\label{ex:p-system-riemann-invariants}
现在考虑 $p$ 方程组
\begin{equation}\label{eq:p-system-vector-form}
\begin{pmatrix}w\\v\end{pmatrix}_t+\begin{pmatrix}-v\\p(w)\end{pmatrix}_x=\begin{pmatrix}0\\0\end{pmatrix}.
\end{equation}
这里有
\begin{equation}\label{eq:p-system-flux-jacobian}
\nabla\bm{f}(w,v)=\begin{pmatrix}0&-1\\p'(w)&0\end{pmatrix}.
\end{equation}
为保证严格双曲性, 假设
\begin{equation}\label{eq:p-system-strict-hyperbolicity}
p'<0.
\end{equation}
此时特征值为 $\lambda_1(w):=-\sqrt{-p'(w)}$ 和 $\lambda_2(w):=\sqrt{-p'(w)}$, 对应的右特征向量为
\begin{equation}\label{eq:p-system-first-eigenvector}
\bm{r}_1(w):=\begin{pmatrix}1\\\sqrt{-p'(w)}\end{pmatrix},
\end{equation}

\begin{equation}\label{eq:p-system-second-eigenvector}
\bm{r}_2(w):=\begin{pmatrix}1\\-\sqrt{-p'(w)}\end{pmatrix}.
\end{equation}
如上述引理所述, 每个特征值对应一个 Riemann 不变量; 它们如下:
\begin{equation}\label{eq:p-system-first-riemann-invariant}
\rho_1(w,v):=v-\Psi(w),
\end{equation}

\begin{equation}\label{eq:p-system-second-riemann-invariant}
\rho_2(w,v):=v+\Psi(w),
\end{equation}
其中
\begin{equation}\label{eq:p-system-riemann-primitive}
\Psi(w):=\int^w\sqrt{-p'(\xi)}\,\mathrm{d}\xi.
\end{equation}

\end{Example}

以下结果给出这一方程组的 Riemann 不变量与特征曲线之间的关系.

\begin{Theorem}
\label{thm:p-system-characteristic-invariants}
设 $(w(x,t),v(x,t))$ 为上述方程组的一个 $C^1$ 解. 那么 Riemann 不变量 $\rho_i(w(x,t),v(x,t))$ 沿满足 $\hat x'(t)=-\lambda_i(w(\hat x(t),t))$ 的特征曲线为常数.
\end{Theorem}

\begin{Proof}
只对 $\rho_1$ 作计算.
\[
\begin{aligned}
\dfrac{\mathrm{d}}{\mathrm{d}t}\rho_1(w(\hat x(t),t),v(\hat x(t),t))
&=v_x\hat x'+v_t-\sqrt{-p'(w)}(w_x\hat x'+w_t)\\
&=-\sqrt{-p'(w)}(w_t-v_x)+(v_t+p'(w)w_x)\\
&=0.
\end{aligned}
\]
对 $\rho_2$ 的计算相同.
\end{Proof}

$p$ 方程组的一个好处是, 可用 Riemann 不变量方便地进行状态空间中的坐标变换; 即由于方程组严格双曲, 有 $\Psi'(w)=\sqrt{-p'(w)}>0$; 因此 $\Psi$ 和映射 $(w,v)\to(\rho_1,\rho_2)$ 都可逆. 若用 $\rho_1,\rho_2$ 将 $p$ 方程组重写, 得到对角方程组
\begin{equation}\label{eq:p-system-diagonal-first}
\rho_{1,t}+\hat\lambda(\rho_1-\rho_2)\rho_{1,x}=0,
\end{equation}

\begin{equation}\label{eq:p-system-diagonal-second}
\rho_{2,t}-\hat\lambda(\rho_1-\rho_2)\rho_{2,x}=0.
\end{equation}
这里
\begin{equation}\label{eq:p-system-diagonal-speed}
\hat\lambda(s):=\sqrt{-p'\left(\Psi^{-1}\left(-\dfrac{s}{2}\right)\right)}.
\end{equation}


Theorem~\ref{thm:p-system-characteristic-invariants} 和上述对角化过程都可推广到任意由两个守恒律组成的严格双曲方程组. 还可用 Riemann 不变量描述一个假设, 物理中产生的守恒律方程组常常满足这个假设.

\begin{Definition}
\label{def:genuine-nonlinearity}
对区域 $D\subset\mathbb{R}^n$, 若守恒律方程组 \eqref{eq:conservation-law-system} 在 $D$ 内满足
\begin{equation}\label{eq:genuine-nonlinearity-definition}
\nabla\lambda_k\cdot\bm{r}_k\ne0,\quad k=1,2,\ldots,n
\end{equation}
则称它在该区域内是\textbf{真正非线性的}.
\end{Definition}

\begin{Example}
\label{ex:scalar-genuine-nonlinearity}
对单个守恒律 \eqref{eq:scalar-conservation-law-quasilinear}, 有 $\lambda(u)=f'(u)$ 和 $r=1$, 所以 $\nabla\lambda(u)\cdot r=f''(u)$. 将满足 $f''>0$ ($<0$) 的函数称为\emph{强凸} (\emph{强凹}) 函数. 在守恒律中, 有时将这样的函数称为严格凸 (严格凹) 函数. (严格地说) 这是不正确的. 因此, $f$ 的强凸性或强凹性都蕴含真正非线性. 出于物理原因, 常假设强凸性. (描述守恒律稳态的变分问题通常表述为极小化问题, 而非极大化问题.)
\end{Example}

\begin{Example}
\label{ex:p-system-genuine-nonlinearity}
对 $p$ 方程组, 有
\begin{equation}\label{eq:p-system-genuine-nonlinearity}
\nabla\lambda_1\cdot\bm{r}_1=-\nabla\lambda_2\cdot\bm{r}_2=\dfrac{p''}{2\sqrt{-p'}}.
\end{equation}
同样, $p$ 的强凸性或强凹性足以保证真正非线性. 在气体动力学的典型应用中, 假设 $p$ 强凸.
\end{Example}

应当指出, 有些有趣的物理问题并不是真正非线性的. 特别地, 在 $p$ 方程组中, 有时假设函数 $p$ 有拐点. 本书不详细讨论这样的问题, 但应向读者介绍以下术语.

\begin{Definition}
\label{def:linear-degeneracy}
若
\begin{equation}\label{eq:linear-degeneracy-definition}
\nabla\lambda_k(\bm{u})\cdot\bm{r}_k(\bm{u})=0.
\end{equation}
则称第 $k$ 个特征场在 $\bm{u}$ 处是\textbf{线性退化的}.
\end{Definition}
